% LaTeX Artifact Appendix — HierRTLBench
% MLCAD 2026
%
% Based on the AE template by Grigori Fursin et al.
% CC BY 4.0 license
%
% Build:
%   pdflatex artifact_appendix.tex
%   pdflatex artifact_appendix.tex   % (second pass for hyperref)

\documentclass[sigplan,10pt,twocolumn]{acmart}

% ── Packages ──────────────────────────────────────────────────────────────────
\usepackage{hyperref}
\usepackage{booktabs}
\usepackage{listings}
\usepackage{xcolor}
\usepackage{microtype}
\usepackage{enumitem}

% ── Listing style for shell / code snippets ───────────────────────────────────
\lstdefinestyle{shell}{
  basicstyle=\ttfamily\scriptsize,
  breaklines=true,
  breakatwhitespace=false,
  columns=flexible,
  keepspaces=true,
  backgroundcolor=\color{gray!8},
  frame=single,
  framesep=2pt,
  rulecolor=\color{gray!40},
  xleftmargin=4pt,
  xrightmargin=4pt,
  aboveskip=3pt,
  belowskip=3pt,
}
\lstset{style=shell}

% ── Suppress ACM copyright / conference info for standalone appendix PDF ──────
\setcopyright{none}
\acmConference{}{}{}
\acmDOI{}
\acmISBN{}

% ── Tighten list spacing ──────────────────────────────────────────────────────
\setlist[itemize]{leftmargin=*, topsep=1pt, itemsep=0pt, parsep=0pt}
\setlist[enumerate]{leftmargin=*, topsep=1pt, itemsep=0pt, parsep=0pt}

% ── Tighten table row spacing and float gaps ──────────────────────────────────
\renewcommand{\arraystretch}{0.92}
\setlength{\textfloatsep}{6pt plus 1pt minus 1pt}
\setlength{\floatsep}{6pt plus 1pt minus 1pt}
\setlength{\intextsep}{6pt plus 1pt minus 1pt}
\setlength{\abovecaptionskip}{2pt}
\setlength{\belowcaptionskip}{2pt}

% ─────────────────────────────────────────────────────────────────────────────
\begin{document}
\footnotesize
\sloppy

% ─────────────────────────────────────────────────────────────────────────────
\appendix
\section{Artifact Appendix}

% ─────────────────────────────────────────────────────────────────────────────
\subsection{Abstract}

This artifact provides the benchmark suite, generation scripts, pre-generated
Verilog outputs, and testbenches for \textit{HierRTLBench}. It compares two
RTL generation strategies, \textit{unified} (one prompt per module) vs.\
\textit{hierarchical} (one prompt per submodule) across 14 hardware
designs, using three open-weight VeriGen models (2B/6B/16B). Reviewers can
(1)~\textbf{validate} pre-generated outputs with \texttt{iverilog} syntax
checks, no GPU needed ($\sim$6~MB disk, Python~3.9+), or
(2)~\textbf{regenerate} outputs with the provided scripts (NVIDIA GPU,
$\geq$16~GB VRAM, $\geq$80~GB for 16B; $\sim$60~GB disk for weights).

% ─────────────────────────────────────────────────────────────────────────────
\subsection{Artifact Checklist (Meta-information)}

\begin{itemize}
  \item \textbf{Algorithm:} Hierarchical LLM prompt decomposition for
        Verilog RTL generation.
  \item \textbf{Program:} 14 modules (ALU, regfile, UART~TX, multi-cycle CPU,
        control unit, pipeline CPU, data cache, RR arbiter, AES-128, SHA-256,
        CRC-32, FFT-8, MatMul-16, bubble sort), each hierarchically
        decomposed into 2--10 submodules (50 submodules total); all
        included.
  \item \textbf{Compilation:} Python~3.9+, PyTorch~2.1+ (CUDA~12.x); see
        \texttt{requirements.txt}. A \texttt{Dockerfile} is provided for the
        cached/CPU-only evaluation path.
  \item \textbf{Transformations:} None. \quad \textbf{Binary:} None included.
  \item \textbf{Model:} VeriGen (CodeGen fine-tuned on Verilog),
        \texttt{shailja/fine-tuned-\allowbreak codegen-\{2B,6B,16B\}-Verilog}
        on Hugging~Face; open-weight, no API key. Auto-downloaded by
        \texttt{transformers} (set \texttt{HF\_HOME} on HPC).
  \item \textbf{LLM prompts/settings:} 50 hierarchical submodule prompts
        (\texttt{verigen\_prompts.txt}); 14 unified prompts inline.
        \texttt{max\_length=2048}, greedy decoding
        (\texttt{temperature=0}, \texttt{do\_sample=False}), 1 sample/submodule.
  \item \textbf{Data set:} 14 module specs embedded as prompt strings; no
        external download.
  \item \textbf{Run-time environment:} Linux (Ubuntu~20.04/22.04); Python~3.9+,
        PyTorch~2.1+, \texttt{transformers}~4.38+, \texttt{accelerate}~0.27+,
        \texttt{sentencepiece}~0.1.99+. \texttt{Dockerfile} provided for
        cached-output validation (CPU only); \texttt{setup\_env.sh} for a
        native venv (needed for GPU regeneration).
  \item \textbf{Hardware:} GPU only for regeneration ($\geq$16~GB VRAM
        2B/6B, $\geq$80~GB 16B; paper used H100~80~GB/SLURM). Syntax
        evaluation runs on any CPU.
  \item \textbf{Run-time state:} Generation uses greedy decoding
        (deterministic); observed $\pm$1 file ($\pm$7\%) variation across
        reruns stems from library/driver version drift (e.g.\
        \texttt{iverilog} version, PyTorch/CUDA build), not sampling.
  \item \textbf{Execution time:} $\sim$1.5/2.5/4~hrs (2B/6B/16B) on one H100;
        syntax evaluation $<$1~min.
  \item \textbf{Metrics:} (1)~Syntax pass rate: \% of hierarchical submodule
        \texttt{.v} files compiling under \texttt{iverilog -tnull}
        (denominator = 50 submodule files per model, not 14 modules), per
        model and module --- reproduced by \texttt{evaluate\_syntax.py}.
        (2)~End-to-end functional pass rate (paper Table~4): \% of top-level
        modules whose integrated design passes testbench simulation with
        \texttt{vvp}; this artifact ships pre-generated outputs and
        testbenches but the integration + simulation step is manual (see
        \S\ref{sec:eval}).
  \item \textbf{Output:} Per-module \texttt{.v} files under
        \texttt{verigen\_out/} (hier.) and \texttt{unified\_outputs/}
        (unified); console summary from \texttt{evaluate\_syntax.py}.
  \item \textbf{Expected results:} 
  \item \textbf{Experiments:} Cached-output validation (no rerun needed) and
        full-rerun path both supported.
  \item \textbf{Proprietary EDA/PDKs:} Synopsys VCS+Verdi used
        \emph{optionally} for coverage only (paper); not required for key
        results. 
  \item \textbf{Disk space:} $\sim$6~MB (repo+outputs); $\sim$50~GB (weights,
        on demand).
  \item \textbf{Workflow prep time:} 10--15~min. \quad
        \textbf{Experiment time:} $<$5~min cached; 6--8 GPU-hrs full regenration.
  \item \textbf{API/GPU cost:} None (open-weight, local); $\sim$6--8
  \item \textbf{Publicly available:} Yes. \quad \textbf{Code license:} MIT License
  \item \textbf{Data/model license:} VeriGen weights follow the Apache 2.0 License
  \item \textbf{Workflow framework:} Python scripts + SLURM job arrays.
  \item \textbf{Zenodo DOI:} (To be added after archiving.)
\end{itemize}

% ─────────────────────────────────────────────────────────────────────────────
\subsection{Description}

\noindent\textbf{Access:}
\begin{lstlisting}
git clone https://github.com/<yashwant>/HierRTLBench.git
\end{lstlisting}
Archived snapshot: Zenodo DOI \texttt{(to be added after archiving)}.
Outputs are pre-included: \texttt{vgen\_project/verigen\_out/\{2B,6B,16B\}/}
(hier.), \texttt{vgen\_project/unified\_outputs/\{2B,6B,16B\}/} (unified),
and \texttt{vgen\_project/verigen\_testbenches/} (14 modules). $\sim$6~MB
after clone; model weights ($\sim$50~GB) download on demand.

\noindent\textbf{Hardware:} CPU only for syntax validation. For
regeneration: NVIDIA GPU, CUDA~12.x, $\geq$16~GB VRAM (2B), $\geq$24~GB
(6B), $\geq$80~GB (16B, H100/A100 recommended).

\noindent\textbf{Software:} Python~3.9+, PyTorch~2.1+, \texttt{transformers}
4.38+, \texttt{accelerate}~0.27+, \texttt{sentencepiece}~0.1.99+,
\texttt{iverilog}~11.0+. Synopsys VCS+Verdi optional (coverage only, not
required for key results; no PDKs used).

\noindent\textbf{Data sets:} 14 module specs embedded in
\texttt{verigen\_runner.py} (50 hier.\ submodule prompts) and
\texttt{run\_unified\_verigen.py} (14 unified prompts); listed in
\texttt{verigen\_prompts.txt}. No external download.

\noindent\textbf{Models:} Three open-weight VeriGen checkpoints on
Hugging~Face (auto-downloaded; set \texttt{HF\_HOME} to redirect cache).
Pinned to the latest revision on the \texttt{main} branch of each
repository at time of writing; no separate commit hash is published
upstream beyond the HF \texttt{main} ref.

\begin{table}[h]
\centering\footnotesize
\begin{tabular}{@{}lll@{}}
\toprule
Model & HF Identifier & Size \\
\midrule
VeriGen-2B  & \texttt{shailja/fine-tuned-codegen-2B-Verilog}  & $\sim$4~GB \\
VeriGen-6B  & \texttt{shailja/fine-tuned-codegen-6B-Verilog}  & $\sim$12~GB \\
VeriGen-16B & \texttt{shailja/fine-tuned-codegen-16B-Verilog} & $\sim$32~GB \\
\bottomrule
\end{tabular}
\end{table}
\vspace{-6pt}
Inference settings (all experiments): \texttt{max\_length=2048}, greedy
decoding (\texttt{temperature=0}, \texttt{do\_sample=False}),
1 sample/submodule, no stop sequences.

% ─────────────────────────────────────────────────────────────────────────────
\subsection{Installation}

\noindent\textbf{Option 1: Docker (recommended for cached validation,
no GPU):}
\begin{lstlisting}
docker build -t hierrtlbench-ae .
# Smoke test (regfile_mem.v is a documented PASS case)
docker run --rm hierrtlbench-ae \
  iverilog -tnull vgen_project/verigen_out/2B/m02_regfile/regfile_mem.v \
  && echo PASS || echo FAIL
# Full cached-output evaluation
docker run --rm hierrtlbench-ae bash run_evaluation.sh --cached
\end{lstlisting}
Expected: \texttt{PASS}; then a pass-rate summary table matching
Table~\ref{tab:expected} within tolerance.

\noindent\textbf{Option 2: Native install (required for GPU
regeneration):}
\begin{lstlisting}
# 1. System deps (Ubuntu/Debian; or brew on macOS)
sudo apt-get install iverilog python3 python3-pip
# 2. Python deps
pip install -r requirements.txt
pip install torch --index-url https://download.pytorch.org/whl/cu121
# 3. HPC only
bash vgen_project/setup_env.sh
# Smoke test (regfile_mem.v is a documented PASS case)
iverilog -tnull vgen_project/verigen_out/2B/m02_regfile/regfile_mem.v \
  && echo PASS || echo FAIL
python evaluate_syntax.py --models 2B
\end{lstlisting}
Expected: \texttt{PASS}; a pass-rate summary table.

\noindent\textbf{Note:} \texttt{m01\_alu/alu\_addsub.v} is \emph{not} a
valid standalone smoke-test target, it is a documented F2 failure
case (references undefined \texttt{half\_adder}/\texttt{full\_adder}
submodules; see paper Section~6) and will not compile in isolation.

% ─────────────────────────────────────────────────────────────────────────────
\subsection{Experiment Workflow}

Three validation paths, all wrapped by \texttt{run\_evaluation.sh}:

\begin{table}[h]
\centering\footnotesize
\begin{tabular}{@{}p{1.55cm}p{2.6cm}p{2.9cm}@{}}
\toprule
Path & Command & Notes \\
\midrule
A: Cached (recommended) & \texttt{run\_evaluation.sh --cached} &
  No GPU; uses included outputs \\
B: Full regen & \texttt{run\_evaluation.sh --full-regen} &
  All 3 models, GPU req'd, 6--8~hrs \\
C: Bounded subset & \texttt{verigen\_runner.py --model 2B --idx 0} &
  Single submodule smoke test \\
\bottomrule
\end{tabular}
\end{table}
\vspace{-6pt}

Path A directly reproduces the paper's key table:
\begin{lstlisting}
python evaluate_syntax.py --hier_dir vgen_project/verigen_out \
  --unified_dir vgen_project/unified_outputs \
  --tb_dir vgen_project/verigen_testbenches/verigen_testbenches
\end{lstlisting}
Path B reruns generation per model
(\texttt{verigen\_runner.py}/\texttt{run\_unified\_verigen.py}
\texttt{--model \{2B,6B,16B\}}), writes integration tops
(\texttt{write\_integrations.py}), then re-evaluates with the same script.
Path C runs one submodule for a fast sanity check (used by SLURM array
jobs via \texttt{--idx}; valid range depends on script version. Run
\texttt{verigen\_runner.py --list} to enumerate).

\noindent\textbf{Paper result} $\rightarrow$ \textbf{script}: hierarchical/
unified pass rates and per-module breakdown all come from
\texttt{evaluate\_syntax.py}; generated Verilog is under
\texttt{verigen\_out/} and \texttt{unified\_outputs/}.

% ─────────────────────────────────────────────────────────────────────────────
\subsection{Evaluation and Expected Results}
\label{sec:eval}

Run Path~A to reproduce the key table. \texttt{evaluate\_syntax.py} checks
each hierarchical \emph{submodule file} independently
(denominator~=~number of generated \texttt{.v} files per model, not~14):

\begin{table}[h]
\centering\footnotesize
\caption{Expected syntax pass rates (hierarchical, per-submodule-file).}
\label{tab:expected}
\begin{tabular}{@{}llll@{}}
\toprule
Model & Passed &  Rate \\
\midrule
VeriGen-2B  &  \\
VeriGen-6B  &      \\
VeriGen-16B & \\
\bottomrule
\end{tabular}
\end{table}
\vspace{-6pt}
Exact pass/total counts depend on the submodule-count reconciliation noted
above (50 per paper vs.\ 64 currently shipped); rates are stable either
way within tolerance.

\noindent\textbf{Tolerance:} 

\noindent\textbf{Functional pass rate (paper Table~4, not scripted):} the
paper additionally reports end-to-end functional pass rates of 30\%~(2B),
0\%~(6B), 28\%~(16B), obtained by integrating each module's submodules
(\texttt{write\_integrations.py}) and simulating the result against the
testbenches in \texttt{vgen\_project/verigen\_testbenches/} with
\texttt{iverilog}/\texttt{vvp} (or Vivado). This integration + simulation
step is not wrapped by \texttt{run\_evaluation.sh} and must be run
per-module manually; reviewers seeking only the syntax-pass-rate headline
result can rely on Path~A alone.

% ─────────────────────────────────────────────────────────────────────────────
\subsection{Experiment Customization}

\begin{itemize}
  \item \textbf{Model size:} \texttt{--model 2B/6B/16B} on any runner.
  \item \textbf{Single submodule:} \texttt{--idx N} in
        \texttt{verigen\_runner.py} (used by SLURM array jobs); run
        \texttt{--list} for the valid range.
  \item \textbf{Output dir:} \texttt{--output\_dir} on all scripts.
  \item \textbf{Temperature:} edit \texttt{GEN\_CFG["temperature"]} in the
        runner scripts.
  \item \textbf{Custom prompts:} extend \texttt{SUBMODULES} in
        \texttt{verigen\_runner.py}.
\end{itemize}

\noindent\textbf{Notes:} SLURM scripts use HPC-specific paths (e.g.\
\texttt{/scratch/soumyaj/}) --- edit \texttt{PYTHON}/\texttt{HF\_HOME} before
running elsewhere.

\end{document}